\chapter{Установка и настройка iRedMail}
\label{ch:practice}

\section{Скачивание и распаковка iRedMail}

На ВМ \texttt{mail} скачиваем архив iRedMail версии 1.6.2:
\begin{lstlisting}
wget https://github.com/iredmail/iRedMail/archive/refs/tags/1.6.2.tar.gz
\end{lstlisting}

Распаковываем:
\begin{lstlisting}
tar xvf 1.6.2.tar.gz
cd iRedMail-1.6.2
\end{lstlisting}

% --- Скриншот: завершение wget (100% downloaded) ---
\begin{figure}[h]
  \centering
  \includegraphics[width=0.85\textwidth]{img/06_wget_iredmail}
  \caption{Скачивание архива iRedMail}
  \label{fig:wget_iredmail}
\end{figure}

\section{Скачивание пакетов iRedMail}

Перед запуском установщика скачиваем зависимости:
\begin{lstlisting}
cd ./pkgs/
chmod +x get_all.sh
./get_all.sh
\end{lstlisting}

% --- Скриншот: завершение get_all.sh ---
\begin{figure}[h]
  \centering
  \includegraphics[width=0.85\textwidth]{img/07_get_all_done}
  \caption{Завершение скачивания пакетов \texttt{pkgs/get\_all.sh}}
  \label{fig:get_all}
\end{figure}

\section{Запуск интерактивного установщика}

Возвращаемся в корневой каталог iRedMail и запускаем установщик:
\begin{lstlisting}
cd ..
chmod +x iRedMail.sh
./iRedMail.sh
\end{lstlisting}

Установщик задаёт ряд вопросов. Ответы для данного варианта:

\subsection{Путь хранения почтовых ящиков}

Оставляем значение по умолчанию:
\begin{lstlisting}
/var/vmail
\end{lstlisting}

% --- Скриншот: экран Default mail storage path ---
\begin{figure}[h]
  \centering
  \includegraphics[width=0.85\textwidth]{img/08_install_storage}
  \caption{Установщик iRedMail --- путь хранения почты}
  \label{fig:install_storage}
\end{figure}

\subsection{Выбор веб-сервера}

Выбираем \textbf{Nginx}:

% --- Скриншот: экран выбора веб-сервера (Nginx выделен) ---
\begin{figure}[h]
  \centering
  \includegraphics[width=0.85\textwidth]{img/09_install_webserver}
  \caption{Установщик iRedMail --- выбор веб-сервера Nginx}
  \label{fig:install_web}
\end{figure}

\subsection{Выбор бэкенда хранения учётных записей}

Выбираем \textbf{OpenLDAP}:

% --- Скриншот: экран выбора бэкенда (OpenLDAP выделен) ---
\begin{figure}[h]
  \centering
  \includegraphics[width=0.85\textwidth]{img/10_install_backend}
  \caption{Установщик iRedMail --- выбор бэкенда OpenLDAP}
  \label{fig:install_backend}
\end{figure}

\subsection{LDAP suffix (root DN)}

Указываем суффикс дерева каталога:
\begin{lstlisting}
dc=yazikov,dc=iks531,dc=local
\end{lstlisting}

% --- Скриншот: экран LDAP suffix с введённым значением ---
\begin{figure}[h]
  \centering
  \includegraphics[width=0.85\textwidth]{img/11_install_ldap_suffix}
  \caption{Установщик iRedMail --- LDAP suffix}
  \label{fig:install_ldap}
\end{figure}

\subsection{Первый почтовый домен}

Указываем домен. Важно: он \textbf{не должен} совпадать с FQDN хоста.
\begin{lstlisting}
yazikov.iks531.local
\end{lstlisting}

% --- Скриншот: экран ввода первого почтового домена ---
\begin{figure}[h]
  \centering
  \includegraphics[width=0.85\textwidth]{img/12_install_mail_domain}
  \caption{Установщик iRedMail --- первый почтовый домен}
  \label{fig:install_domain}
\end{figure}

\subsection{Опциональные компоненты}

Выбираем: \textbf{Roundcubemail}, \textbf{netdata}, \textbf{iRedAdmin},
\textbf{Fail2ban}. SOGo оставляем невыбранным.

% --- Скриншот: экран Optional components ---
\begin{figure}[h]
  \centering
  \includegraphics[width=0.85\textwidth]{img/13_install_components}
  \caption{Установщик iRedMail --- выбор компонентов}
  \label{fig:install_components}
\end{figure}

\subsection{Подтверждение и завершение}

Подтверждаем установку (\texttt{y}) и разрешаем настройку firewall (\texttt{y}).
Установка занимает 10--20 минут в зависимости от скорости интернета.
По завершении установщик запрашивает перезагрузку --- соглашаемся.

% --- Скриншот: финальный экран установщика (Installation Completed) ---
\begin{figure}[h]
  \centering
  \includegraphics[width=0.85\textwidth]{img/14_install_done}
  \caption{iRedMail --- установка завершена успешно}
  \label{fig:install_done}
\end{figure}

\section{Проверка сервисов после перезагрузки}

После перезагрузки ВМ \texttt{mail} проверяем статус ключевых сервисов:
\begin{lstlisting}
systemctl status postfix
systemctl status dovecot
systemctl status nginx
systemctl status slapd
\end{lstlisting}

% --- Скриншот: systemctl status postfix (active running) ---
\begin{figure}[h]
  \centering
  \includegraphics[width=0.85\textwidth]{img/15_services_status}
  \caption{Статус сервисов после установки iRedMail}
  \label{fig:svc_status}
\end{figure}

\section{Вход в панель iRedAdmin}

С ВМ \texttt{Desktop} открываем браузер Firefox и переходим по адресу:
\begin{lstlisting}
https://mail.yazikov.iks531.local/iredadmin
\end{lstlisting}

Браузер выдаёт предупреждение о самоподписанном сертификате ---
нажимаем \textit{Advanced} → \textit{Accept the Risk and Continue}.

Входим:
\begin{itemize}
  \item Логин: \texttt{postmaster@yazikov.iks531.local}
  \item Пароль: указанный при установке
\end{itemize}

% --- Скриншот: форма входа iRedAdmin в браузере ---
\begin{figure}[h]
  \centering
  \includegraphics[width=0.85\textwidth]{img/16_iredadmin_login}
  \caption{Панель iRedAdmin --- форма входа}
  \label{fig:iredadmin_login}
\end{figure}

\section{Создание почтового пользователя}

В iRedAdmin переходим: \textit{Users} → \textit{Add user}.
Заполняем форму:
\begin{itemize}
  \item Email: \texttt{user1@yazikov.iks531.local}
  \item Пароль: задаём произвольный
\end{itemize}
Нажимаем \textit{Add}.

% --- Скриншот: форма создания пользователя в iRedAdmin ---
\begin{figure}[h]
  \centering
  \includegraphics[width=0.85\textwidth]{img/17_iredadmin_add_user}
  \caption{iRedAdmin --- создание нового пользователя}
  \label{fig:add_user}
\end{figure}

% --- Скриншот: список пользователей с созданным user1 ---
\begin{figure}[h]
  \centering
  \includegraphics[width=0.85\textwidth]{img/18_iredadmin_user_list}
  \caption{iRedAdmin --- список пользователей домена}
  \label{fig:user_list}
\end{figure}

\section{Отправка тестового письма через Roundcube}

Переходим в веб-почту:
\begin{lstlisting}
https://mail.yazikov.iks531.local/mail
\end{lstlisting}

Входим как \texttt{postmaster@yazikov.iks531.local}.

% --- Скриншот: интерфейс Roundcubemail после входа ---
\begin{figure}[h]
  \centering
  \includegraphics[width=0.85\textwidth]{img/19_roundcube_inbox}
  \caption{Roundcubemail --- входящие письма}
  \label{fig:roundcube_inbox}
\end{figure}

Нажимаем \textit{Compose}, заполняем:
\begin{itemize}
  \item To: \texttt{user1@yazikov.iks531.local}
  \item Subject: \texttt{Тест ЛР-7}
  \item Тело письма: любой текст
\end{itemize}
Отправляем.

% --- Скриншот: окно создания письма в Roundcube ---
\begin{figure}[h]
  \centering
  \includegraphics[width=0.85\textwidth]{img/20_roundcube_compose}
  \caption{Roundcubemail --- окно создания письма}
  \label{fig:roundcube_compose}
\end{figure}

\section{Получение письма}

Выходим из аккаунта \texttt{postmaster} и входим как
\texttt{user1@yazikov.iks531.local}.
Открываем Inbox --- видим пришедшее письмо.

% --- Скриншот: Inbox user1 с полученным письмом ---
\begin{figure}[h]
  \centering
  \includegraphics[width=0.85\textwidth]{img/21_roundcube_received}
  \caption{Roundcubemail --- полученное письмо в Inbox \texttt{user1}}
  \label{fig:roundcube_received}
\end{figure}

% --- Скриншот: открытое письмо (тема и содержимое) ---
\begin{figure}[h]
  \centering
  \includegraphics[width=0.85\textwidth]{img/22_roundcube_open_letter}
  \caption{Roundcubemail --- открытое письмо}
  \label{fig:roundcube_open}
\end{figure}

\section{Просмотр почтового лога}

Для проверки прохождения письма через Postfix смотрим лог на ВМ \texttt{mail}:
\begin{lstlisting}
tail -30 /var/log/mail.log
\end{lstlisting}

% --- Скриншот: tail /var/log/mail.log (статус=sent) ---
\begin{figure}[h]
  \centering
  \includegraphics[width=0.85\textwidth]{img/23_mail_log}
  \caption{Лог Postfix --- успешная доставка письма (\texttt{status=sent})}
  \label{fig:mail_log}
\end{figure}

\endinput
